\documentclass[12pt,a4paper]{article}

\usepackage[spanish,es-nodecimaldot]{babel}
\usepackage[utf8]{inputenc}
\usepackage[T1]{fontenc}
\usepackage{lmodern}
\usepackage{geometry}
\usepackage{microtype}
\usepackage{setspace}
\usepackage{parskip}
\usepackage{enumitem}
\usepackage{xcolor}
\usepackage{listings}
\usepackage{hyperref}
\usepackage{amsmath}
\usepackage{amssymb}
\usepackage{array}
\usepackage{booktabs}
\usepackage{longtable}
\usepackage{tcolorbox}

\geometry{margin=2.5cm}
\onehalfspacing
\setlist[itemize]{leftmargin=*,itemsep=0.2em}
\setlist[enumerate]{leftmargin=*,itemsep=0.2em}

\hypersetup{
  colorlinks=true,
  linkcolor=black,
  urlcolor=blue,
  pdftitle={Resumen teorico TP6 - Testing basado en sintaxis y mutacion},
  pdfauthor={}
}

\newtcolorbox{definicionbox}[1][]{
  colback=blue!4!white,
  colframe=blue!50!black,
  fonttitle=\bfseries,
  title=Definicion,
  #1
}

\newtcolorbox{ideabox}[1][]{
  colback=green!4!white,
  colframe=green!40!black,
  fonttitle=\bfseries,
  title=Idea clave,
  #1
}

\newtcolorbox{ejemplobox}[1][]{
  colback=orange!4!white,
  colframe=orange!50!black,
  fonttitle=\bfseries,
  title=Ejemplo,
  #1
}

\lstdefinestyle{codestyle}{
  basicstyle=\ttfamily\small,
  frame=single,
  breaklines=true,
  showstringspaces=false
}

\title{\textbf{Resumen Teorico -- Practico 6}\\
Testing basado en sintaxis, mutacion y fuzzing\\
\textit{Testing de Software 2024}}
\author{}
\date{}

\begin{document}

\maketitle
\tableofcontents
\newpage

\section*{Introduccion}
\addcontentsline{toc}{section}{Introduccion}

Este resumen sintetiza el material teorico asociado al TP6:

\begin{itemize}
  \item \textit{Introduction to Software Testing} (Ammann \& Offutt, 2da edicion), Capitulo 9:
  \textbf{Syntax-Based Testing}.
  \item \textit{Notas 10: Testing basado en sintaxis} (V. Bengolea).
  \item \textit{Notas 11: Testing basado en mutacion} (V. Bengolea).
\end{itemize}

\textbf{Idea general.} Muchos artefactos del software (programas, inputs, mensajes de
protocolo, consultas SQL, archivos de configuracion) responden a una \emph{sintaxis estricta},
descripta tipicamente por una gramatica BNF. El \textbf{testing basado en sintaxis} usa esa
estructura como modelo y deriva tests de dos formas complementarias:

\begin{itemize}
  \item \textbf{Cubriendo} la gramatica: generar cadenas que ejerciten todas las producciones,
  todos los simbolos terminales, todos los pares de producciones, etc.
  \item \textbf{Violando} la gramatica: mutarla o mutar las cadenas para generar tests
  \emph{invalidos} que sirven como prueba de robustez.
\end{itemize}

La aplicacion mas conocida es el \textbf{testing de mutacion}, donde lo que se muta es el
\emph{propio programa} y la suite debe ser capaz de distinguir el programa original de cada
variante. El TP6 cierra con \textbf{fuzzing}, que es la realizacion practica de generacion
aleatoria/sintactica de inputs.

\section{Testing basado en sintaxis: vista general}

\subsection{Cuatro estructuras para modelar software}

Recordando la columna vertebral del libro, hay cuatro grandes formas de modelar software para
generar tests:

\begin{enumerate}
  \item \textbf{Espacio de inputs} (TP2): particionar el dominio.
  \item \textbf{Grafos} (TP4): CFG, FSM, statecharts.
  \item \textbf{Logica} (TP5): predicados y clausulas.
  \item \textbf{Sintaxis} (TP6): gramaticas y derivaciones.
\end{enumerate}

\begin{ideabox}
La sintaxis se aplica a cuatro fuentes tipicas: \emph{codigo}, \emph{modelos} (UML, diagramas),
\emph{inputs} permitidos por el programa, y \emph{tests de integracion} (mensajes entre
componentes).
\end{ideabox}

\subsection{Por que sintaxis}

\begin{itemize}
  \item Cuando hay una gramatica explicita, la generacion de tests es \emph{mecanica}: se puede
  derivar automaticamente un conjunto enorme de cadenas validas.
  \item Cuando se quiere testear el \emph{manejo de errores}, se mutan las cadenas (o la
  gramatica) y se observa que el sistema rechace o reaccione adecuadamente.
  \item Cuando se quiere medir la \emph{calidad de una suite}, se mutan los programas (testing
  de mutacion) y se pide que la suite distinga el original del mutante.
\end{itemize}

\section{Gramaticas BNF}

\begin{definicionbox}
Una \textbf{gramatica BNF} (Backus-Naur Form) esta dada por:
\begin{itemize}
  \item un conjunto de \emph{simbolos no terminales} $N$,
  \item un conjunto de \emph{simbolos terminales} $T$ disjunto de $N$,
  \item un \emph{simbolo inicial} $S \in N$,
  \item un conjunto de \emph{producciones} de la forma
  $\textit{nt} ::= \alpha_1\ |\ \alpha_2\ |\ \cdots\ |\ \alpha_k$, donde cada $\alpha_i$ es una
  secuencia de simbolos en $N \cup T$.
\end{itemize}
El \emph{lenguaje} de la gramatica es el conjunto de cadenas de terminales obtenibles
derivando desde $S$ por sustituciones sucesivas.
\end{definicionbox}

\textbf{Notacion auxiliar.}

\begin{itemize}
  \item $*$ (estrella de Kleene): cero o mas repeticiones.
  \item $+$: una o mas repeticiones.
  \item $?$ o $[\,]$: opcional.
  \item $|$: alternativa.
\end{itemize}

\begin{ejemplobox}
Gramatica para un stream de acciones:
\begin{lstlisting}[style=codestyle]
Stream ::= accion*
accion ::= accG | accB
accG   ::= "G" s n
accB   ::= "B" t n
s      ::= digito^{1..3}
t      ::= digito^{1..3}
n      ::= digito^2 "." digito^2 "." digito^2
digito ::= "0"|"1"|"2"|"3"|"4"|"5"|"6"|"7"|"8"|"9"
\end{lstlisting}
Una cadena valida es \texttt{G 26 080190 B 22 062794 G 22 112194 B 13 010903}.
\end{ejemplobox}

\subsection{Reconocer vs.\ generar}

Toda gramatica define dos tareas duales:

\begin{itemize}
  \item \textbf{Reconocedor}: dada una cadena, decide si pertenece al lenguaje (\emph{parsear}).
  Existen herramientas (yacc/bison, ANTLR) y los programas las usan para validar inputs.
  \item \textbf{Generador}: dada la gramatica, deriva cadenas que estan en el lenguaje.
\end{itemize}

\section{Criterios de cobertura sobre gramaticas}

\begin{definicionbox}
\textbf{Terminal Symbol Coverage (TSC).} El conjunto de tests $RT$ debe contener, para cada
simbolo terminal $t$ de la gramatica $G$, alguna cadena que lo use.
\end{definicionbox}

\begin{definicionbox}
\textbf{Production Coverage (PDC).} Para cada produccion $p$ de la gramatica, $RT$ debe
contener alguna cadena cuya derivacion use $p$.
\end{definicionbox}

\begin{ideabox}
PDC subsume TSC: para usar todas las producciones es necesario haber usado tarde o temprano
todos los simbolos terminales que aparecen en ellas. PDC es el criterio mas frecuente porque
se mapea directo a las ``reglas distintas'' de la gramatica.
\end{ideabox}

\textbf{Criterios mas exigentes.}

\begin{itemize}
  \item \textbf{Derivation Coverage (DC).} Que cada \emph{derivacion completa} posible aparezca
  en alguna cadena. En general infinita; se usa solo en gramaticas chicas o restringidas.
  \item \textbf{Cobertura de pares} de producciones, analoga a la cobertura de pares de aristas
  en grafos.
\end{itemize}

\subsection{Cadenas mutadas vs.\ gramaticas mutadas}

Cuando la cobertura no apunta a \emph{validar} sino a \emph{detectar manejo erroneo}, se
generan cadenas \emph{invalidas}. Hay dos formas:

\begin{itemize}
  \item \textbf{Mutar la gramatica} (cambiar una regla) y generar a partir de la version
  mutada: las cadenas producidas \emph{no} pertenecen al lenguaje original.
  \item \textbf{Mutar cadenas basicas} (validas), generando variantes que pueden quedar
  validas o invalidas dependiendo del cambio.
\end{itemize}

\begin{ideabox}
El esquema general es: ``\textit{partir de cadenas/derivaciones validas, aplicar mutaciones y
ver si el sistema bajo prueba detecta o tolera el cambio}''. Si las mutaciones siempre
producen cadenas invalidas, el oraculo es ``el programa debe rechazar''. Si pueden quedar
validas, el oraculo es ``el programa debe seguir funcionando bien''.
\end{ideabox}

\section{Mutacion sobre programas}

La aplicacion historica y mas popular del testing basado en sintaxis es \textbf{mutar el
codigo} del programa bajo test.

\begin{definicionbox}
\textbf{Programa mutante.} Variante del programa $P$ obtenida aplicando un \emph{operador de
mutacion} sobre una pequena porcion del codigo. Por ejemplo, cambiar \texttt{\&\&} por
\texttt{||}, reemplazar \texttt{>} por \texttt{>=}, eliminar una asignacion, sustituir una
constante por otra.
\end{definicionbox}

\begin{definicionbox}
\textbf{Matar un mutante.} Dado un mutante $m$ del programa $P$ y un test $t$, se dice que $t$
\emph{mata} a $m$ si la salida de $t$ sobre $P$ difiere de la salida de $t$ sobre $m$.
\end{definicionbox}

\subsection{Tipos de mutantes}

\begin{itemize}
  \item \textbf{Mutante muerto}: hay al menos un test en la suite que lo distingue del
  original.
  \item \textbf{Mutante nacido muerto (\emph{stillborn})}: sintacticamente invalido --- el
  compilador lo rechaza. No aporta al score.
  \item \textbf{Mutante trivial}: casi cualquier test lo mata. Suelen ser baratos pero poco
  utiles porque no exigen mucho a la suite.
  \item \textbf{Mutante equivalente}: tiene exactamente el mismo comportamiento observable que
  el original. \emph{No se puede matar} con ningun test. Hay que identificarlo manualmente y
  descontarlo del denominador del score.
\end{itemize}

\begin{ejemplobox}
En \texttt{int abs(int x) \{ return x < 0 ? -x : x; \}}, mutar \texttt{x < 0} a
\texttt{x <= 0} da un mutante equivalente: la rama \texttt{-x} con $x = 0$ devuelve $0$, igual
que el original.
\end{ejemplobox}

\subsection{Mutation Score}

\begin{definicionbox}
\textbf{Mutation Score} de una suite $T$ sobre un conjunto de mutantes $M$:
\[
\text{score}(T, M) = \frac{|\{ m \in M : T \text{ mata } m \}|}{|M| - |\text{equivalentes}|}
\]
El denominador se ajusta para no contar los mutantes equivalentes (que ningun test puede
matar).
\end{definicionbox}

\begin{ideabox}
Un \emph{mutation score} cercano al $100\%$ indica que la suite es sensible a
\emph{pequenas perturbaciones} del codigo: hay algun input que las descubre. Es un indicador
mas fuerte que cobertura de lineas o aristas, porque obliga al test a \emph{notar el cambio en
la salida}, no solo a ejecutar la linea.
\end{ideabox}

\subsection{Operadores tipicos}

Para programas Java, herramientas como PIT incluyen operadores como:

\begin{itemize}
  \item \textbf{Conditionals Boundary}: $<$ a $\le$, $>$ a $\ge$, etc.
  \item \textbf{Negate Conditionals}: invertir la condicion de un \texttt{if}.
  \item \textbf{Math}: $+ \to -$, $* \to /$, etc.
  \item \textbf{Increments}: $i\texttt{++}$ a $i\texttt{--}$.
  \item \textbf{Return Values}: cambiar \texttt{return x} a \texttt{return 0} o
  \texttt{return null}.
  \item \textbf{Void Method Calls}: eliminar llamadas a metodos que retornan \texttt{void}.
  \item \textbf{Empty Returns}, \textbf{Primitive Returns}, \textbf{Null Returns}, etc.
\end{itemize}

\subsection{Mutacion fuerte vs.\ debil}

\begin{itemize}
  \item \textbf{Strong mutation}: requiere que el cambio se observe \emph{en la salida final}
  del programa. Es el modelo estandar y el que usa PIT.
  \item \textbf{Weak mutation}: basta con que en algun punto intermedio el estado del programa
  difiera del original (mas barato, pero da scores menos significativos).
\end{itemize}

\subsection{Costo y por que vale la pena}

Mutacion es caro: hay que ejecutar la suite una vez por cada mutante. Herramientas modernas
(PIT) mitigan esto con:

\begin{itemize}
  \item \emph{Bytecode mutation}: muta el \texttt{.class} compilado, no el \texttt{.java}.
  \item \emph{Test selection}: solo corre los tests que cubren la linea mutada.
  \item \emph{Incremental analysis}: reusa resultados entre ejecuciones.
\end{itemize}

\section{Fuzzing}

\subsection{Idea general}

\textbf{Fuzzing} es generacion masiva de inputs (validos, invalidos o ambos) para alimentar a
un programa y observar si crashea o se comporta mal. Es una tecnica orientada a robustez y
seguridad mas que a verificacion de correccion funcional.

\subsection{Random fuzzing}

\begin{definicionbox}
Un \textbf{random fuzzer} genera cadenas \emph{sin ninguna estructura previa}: elige una
longitud al azar y rellena con caracteres aleatorios de un alfabeto dado.
\end{definicionbox}

Es facil de implementar pero rara vez logra que la entrada sea sintacticamente valida (para
gramaticas no triviales). Por eso suele explorar mas el manejo de errores que el flujo
``feliz''.

\subsection{Mutation fuzzing}

\begin{definicionbox}
Un \textbf{mutation fuzzer} parte de un conjunto de \emph{semillas} (inputs validos conocidos)
y aplica una secuencia de \emph{mutaciones pequenas} para generar nuevos inputs. Cada mutacion
puede borrar, insertar o modificar caracteres / tokens.
\end{definicionbox}

\textbf{Operadores de mutacion sobre strings} tipicos:

\begin{itemize}
  \item \texttt{deleteRandomCharacter(s)}: elimina un caracter al azar.
  \item \texttt{insertRandomCharacter(s)}: inserta un caracter al azar.
  \item \texttt{flipRandomCharacter(s)}: cambia un caracter (por ejemplo flip de un bit).
\end{itemize}

Aplicar pocas mutaciones a una semilla valida produce inputs ``casi validos'' que ejercitan
mucho mas codigo que strings totalmente aleatorios.

\subsection{Random vs.\ mutation: el tradeoff}

\begin{itemize}
  \item \textbf{Random fuzzing} explora un espacio grande pero acierta poco. Bueno para
  detectar crashes en el lexer/parser.
  \item \textbf{Mutation fuzzing} explora cerca de los inputs validos. Bueno para detectar
  bugs en logica de negocio o partes profundas del codigo.
\end{itemize}

\subsection{El oraculo en fuzzing}

Un fuzzer genera miles de inputs; no se puede definir el resultado correcto para cada uno. El
\emph{oraculo} suele ser:

\begin{itemize}
  \item \textbf{No crash}: el programa no debe terminar con \texttt{SIGABRT} (exit $134$),
  \texttt{SIGSEGV} (exit $139$), ni emitir mensajes de error catastroficos en \texttt{stderr}
  (``segmentation fault'', ``core dumped'', ``illegal instruction'').
  \item \textbf{Sin assertion violations} (en codigo con \texttt{assert}).
  \item \textbf{Sin uso indebido} de memoria (detectable con sanitizers).
\end{itemize}

\begin{ideabox}
Es importante que el oraculo \emph{no} exija que la salida sea ``correcta'' en sentido
funcional: la mayoria de los inputs fuzzed son invalidos, y la respuesta correcta es un error
controlado. Lo que el fuzzer busca es \emph{fallas que el programa no manejo}.
\end{ideabox}

\subsection{Seguridad}

El \emph{warning} del enunciado es real: si se conecta un fuzzer a un programa que puede
\emph{modificar el sistema} (e.g.\ \texttt{rm}, \texttt{dd}, un script con permisos), se corre
el riesgo de que alguna entrada generada provoque dano real. Por eso se eligen programas
\emph{seguros} para fuzzear (un parser, una calculadora como \texttt{bc}, un sandbox, etc.).

\section{Pitest en la practica}

\subsection{Que hace}

\texttt{Pitest} es una herramienta de testing de mutacion para Java. Trabaja sobre el bytecode
compilado, lo cual le da velocidad y soporte para los grandes operadores de mutacion sin
modificar el fuente.

\subsection{Configuracion tipica (\texttt{pom.xml})}

\begin{lstlisting}[style=codestyle]
<plugin>
  <groupId>org.pitest</groupId>
  <artifactId>pitest-maven</artifactId>
  <version>1.8.0</version>
  <configuration>
    <targetClasses><param>my.pkg.MyClass</param></targetClasses>
    <targetTests><param>my.pkg.MyClassTests</param></targetTests>
  </configuration>
  <dependencies>
    <dependency>
      <groupId>org.pitest</groupId>
      <artifactId>pitest-junit5-plugin</artifactId>
      <version>0.16</version>
    </dependency>
  </dependencies>
</plugin>
\end{lstlisting}

\subsection{Ejecucion}

\begin{lstlisting}[style=codestyle]
mvn test org.pitest:pitest-maven:mutationCoverage
\end{lstlisting}

PIT genera un reporte HTML con: total de mutantes, cuantos KILLED, SURVIVED, NO\_COVERAGE,
TIMED\_OUT, MEMORY\_ERROR, etc. Para apuntar a una clase distinta sin tocar el \texttt{pom.xml}
se sobrescribe en linea de comando con \texttt{-DtargetClasses} y \texttt{-DtargetTests}.

\subsection{Como leer el reporte}

\begin{itemize}
  \item \texttt{KILLED}: el test que lo mato esta indicado.
  \item \texttt{SURVIVED}: hay que analizarlo --- o agregar un test que lo distinga, o
  justificar que es equivalente.
  \item \texttt{NO\_COVERAGE}: ningun test ejecuta esa linea; antes de mutacion, conviene
  primero cubrir.
  \item \texttt{TIMED\_OUT}: el mutante introdujo un bucle infinito; cuenta como muerto.
\end{itemize}

\section{Comparacion con criterios anteriores}

\begin{center}
\begin{tabular}{lll}
\toprule
Criterio & Modelo & Pregunta clave \\
\midrule
Particion del input & dominio & \textquestiondown{}cubri todas las clases de equivalencia? \\
NC, EC, EPC, PPC & grafo (CFG) & \textquestiondown{}recorri todos los nodos/aristas? \\
PC, CC, CACC, RACC & predicados logicos & \textquestiondown{}cada clausula importa? \\
TSC, PDC & gramatica & \textquestiondown{}use todas las producciones? \\
Mutation Score & programa mutado & \textquestiondown{}mis tests notan cambios pequenos? \\
Fuzzing & generacion masiva & \textquestiondown{}hay alguna entrada que crashee? \\
\bottomrule
\end{tabular}
\end{center}

\begin{ideabox}
La cobertura de mutacion \emph{subsume} en la practica a casi todos los criterios anteriores:
para matar mutantes hace falta haber recorrido las lineas (NC), las aristas (EC) y haber
hecho que las decisiones tomen distintos valores (PC/CC). No es una equivalencia formal, pero
empiricamente las suites con mutation score alto suelen tener cobertura estructural alta.
\end{ideabox}

\section{Glosario}

\begin{itemize}
  \item \textbf{Operador de mutacion}: regla que produce un mutante a partir del programa
  original cambiando un constructo pequeno.
  \item \textbf{Mutante equivalente}: variante con comportamiento observable identico al
  original.
  \item \textbf{Mutation Score}: fraccion de mutantes muertos, descontando equivalentes.
  \item \textbf{Seed (semilla)}: input valido que se usa como punto de partida del mutation
  fuzzing.
  \item \textbf{Oraculo}: criterio que decide si una corrida es ``correcta'' o ``fallida''. En
  fuzzing suele reducirse a ``no crash''.
  \item \textbf{Stillborn}: mutante sintacticamente invalido.
\end{itemize}

\section*{Conclusion}
\addcontentsline{toc}{section}{Conclusion}

El testing basado en sintaxis es un puente entre la estructura formal de los artefactos
(gramaticas, codigo fuente) y la generacion automatica de tests. La mutacion de programas
proporciona un criterio fuerte de calidad de suite: una suite con \emph{mutation score} alto
no solo recorre el codigo, sino que es \emph{sensible} a cambios. El fuzzing, por su parte,
escala la generacion de inputs a miles de casos y se apoya en oraculos minimos para detectar
fallas graves. Las tres tecnicas son complementarias: la mutacion mide la suite, el fuzzing
estresa el sistema, y la cobertura sintactica garantiza que ningun rincon de la gramatica de
inputs quede sin tocar.

\end{document}
